\section{Metodología}
\label{sec:metodologia}

Este capítulo describe el conjunto de datos, el pipeline de representación, los
parámetros del sistema, la auditoría de fidelidad y las métricas usadas para evaluar el
sistema EAM-TMS de ocho agentes.

% ---------------------------------------------------------------
\subsection{Conjunto de datos}
\label{sec:dataset}

Se utilizó el conjunto de datos ETH-80~\cite{leibe2003eth80} completo, con sus
\textbf{ocho clases}: \texttt{apple}, \texttt{car}, \texttt{cow}, \texttt{cup},
\texttt{dog}, \texttt{horse}, \texttt{pear} y \texttt{tomato}.
Cada clase cuenta con 410 imágenes de $128\times128$ píxeles.
Las imágenes de cada clase se dividieron en 328 imágenes de entrenamiento y 82 de prueba.
La partición se realizó a nivel de archivo de imagen.
Por tanto, no hay solapamiento de imágenes entre train y test, pero no se impone una
separación leave-object-out entre objetos físicos de ETH-80.

Dentro de la fase de alimentación, se definieron 3 particiones disjuntas con propósitos distintos:
\begin{itemize}
    \item \textbf{Llenado} (\textit{filling}): imágenes \texttt{train[0:200]} de cada
    clase, utilizadas para registrar patrones en las memorias de dominio de cada agente
    especialista.
    Cada imagen se registra junto con variantes de \emph{augmentación} (imagen espejo y
    dos rotaciones de $\pm 12^{\circ}$), lo que densifica la distribución latente
    acumulada sin tocar el mecanismo de la memoria: es más percepción, no otro modelo.
    \item \textbf{Interacciones visuales}: imágenes \texttt{train[200:328]} de cada
    clase, usadas para simular las interacciones de percepción durante la formación del
    directorio visual (1024 imágenes en total, 128 por clase).
    \item \textbf{Evaluación}: las 82 imágenes de prueba (\texttt{test}) de cada clase
    (656 en total), nunca vistas durante el llenado ni la formación del directorio.
\end{itemize}

% ---------------------------------------------------------------
\subsection{Pipeline semántico}
\label{sec:pipeline-semantico}

Cada consulta en lenguaje natural transita por un pipeline de cuatro etapas antes de
llegar a las memorias asociativas.

\paragraph{Extracción de conceptos con ConceptNet.}
Se consultó ConceptNet~5.7~\cite{speer2017conceptnet} con las relaciones \texttt{IsA},
\texttt{HasProperty}, \texttt{RelatedTo}, \texttt{HasPart}, \texttt{HasA},
\texttt{MadeOf} y \texttt{CapableOf}.
La extracción de esta versión es \emph{bidireccional}: el concepto de la clase puede
aparecer en cualquiera de los dos extremos de la aserción, lo que multiplica los
candidatos de las clases con vocabulario pobre (\texttt{pear} pasó de 11 a 70 candidatos
crudos; \texttt{tomato} de 12 a 93).
El peso de cada label es su \textbf{masa asociativa} (la suma de los pesos de todas sus
aserciones, no el máximo), y los empates se ordenan por similitud coseno al centroide de
los labels fuertes de la clase.
Los candidatos sin vector fastText real se excluyen \emph{antes} del corte, de modo que
las ocho clases quedan con $\approx$21 labels efectivos y ninguno sintético.

\paragraph{Vectorización con fastText y cuantización por magnitud.}
Cada etiqueta del vocabulario se representó con su vector crudo de 300 dimensiones del
modelo fastText preentrenado~\cite{bojanowski2017fasttext}.
A diferencia de la versión anterior (binarización por signo, que colapsaba cada
componente a 2 de los 16 niveles disponibles), esta versión cuantiza \emph{por
magnitud}: cada vector se escala como $v/S$ con una \textbf{escala global persistida}
$S$ (el percentil 99 de las magnitudes de componente sobre el vocabulario completo,
$S = 0.1864$), se recorta a $[-1, 1]$ y se mapea uniformemente a $m = 16$ niveles.
La misma escala se aplica en el llenado y en la consulta, lo que garantiza que ambos
lados hablen el mismo alfabeto.
Con esto los 16 niveles se pueblan efectivamente (véase la
Figura~\ref{fig:espacio-semantico}).

\paragraph{Procesamiento de consultas con spaCy.}
Las consultas en lenguaje natural se preprocesaron con spaCy~\cite{honnibal2020spacy}:
se aplicó lematización y se retuvieron tokens de categoría gramatical \texttt{NOUN},
\texttt{ADJ} o \texttt{PROPN}, con eliminación de palabras vacías y deduplicación estable.
Se añadió un \emph{rescate de atención léxica}: un token descartado por el etiquetador
POS se rescata si, aislado, se etiqueta como \texttt{NOUN}/\texttt{ADJ}/\texttt{PROPN},
o si es un label registrado en las memorias (este segundo criterio solo \emph{admite}
pistas, nunca censura).
El rescate corrige errores del etiquetador en frases nominales cortas (p.ej.\
\texttt{pear} etiquetado como verbo en ``a bosc pear\ldots''), que de otro modo
eliminarían la pista más distintiva de la consulta antes de llegar a la memoria.

\paragraph{Representación de tokens, sin filtro de vocabulario.}
El vocabulario de ConceptNet sirve para \emph{llenar} las memorias, no para decidir qué
consultas son conocidas.
Cada token de la consulta se convierte en una pista vectorial tomando su vector del
modelo fastText completo, exista o no como etiqueta de ConceptNet.
No se aplica ningún filtro de pertenencia léxica: una vez obtenida la pista, es la
propia memoria la que, mediante el reconocimiento y el \textit{containment}, decide si
hay soporte o se rechaza (puntuación cero).
Solo queda fuera lo que no es representable: un token sin vector fastText real no genera
pista, lo cual es una frontera del codificador lingüístico y no un rechazo de la memoria.

% ---------------------------------------------------------------
\subsection{Percepción visual: encoder ResNet-18}
\label{sec:percepcion}

Para obtener representaciones latentes de las imágenes se entrenó un autoencoder
convolucional con encoder basado en ResNet-18~\cite{he2016resnet} sobre las ocho clases
(no es un \textit{masked autoencoder}: no hay enmascaramiento de la entrada).
El espacio latente tiene dimensión 64.
La pérdida combina reconstrucción y una señal auxiliar de clasificación
($\text{MSE} + 0.1\cdot\text{CE}$); la cabeza clasificadora solo guía el entrenamiento
del codificador y \textbf{no participa} en el enrutamiento ni en la recuperación del
sistema de memorias.
El entrenamiento es \emph{determinista} (semilla 42 en \texttt{torch}, \texttt{numpy} y
\texttt{random}, con cuDNN determinista), de modo que una réplica desde cero es
reproducible; el encoder oficial de esta versión alcanza RMSE 0.110 y precisión de la
cabeza auxiliar de 99.7\%.
Los vectores latentes se cuantizaron a $m = 32$ niveles mediante normalización global
min-max con un margen del 10\%, calculada exclusivamente sobre el pool de llenado
(\texttt{train[0:200]} por clase).
Esta cuantización global evita que la escala relativa entre clases afecte la codificación.

% ---------------------------------------------------------------
\subsection{Banco de evaluación semántica}
\label{sec:banco}

Se construyó un banco de \textbf{411 consultas} en lenguaje natural con etiqueta de
clase correcta conocida (\textit{ground truth}), aproximadamente 50 por dominio,
presentadas de forma intercalada.
Las consultas varían en longitud y vocabulario (p.ej.\ ``a red crunchy fruit'',
``a barking domestic pet'', ``a mug for drinking coffee'') y cubren tanto descripciones
directas como perifrásticas; las clases con léxico compartido (\texttt{pear},
\texttt{tomato}) incluyen consultas con contenido \emph{distintivo}
(``a bosc pear with a narrow neck'', ``tomato on the vine''), reflejando la comunicación
realista de un grupo de Wegner: quien pregunta aporta las señas del dominio.
Para los experimentos de protocolo completo (Secciones~\ref{sec:proto-completo}
a~\ref{sec:formacion}) se usa el subconjunto de las primeras 80 consultas
(10 por clase); el estudio de calibración (Sección~\ref{sec:calibracion}) usa el banco
completo.

% ---------------------------------------------------------------
\subsection{Parámetros de la EAM}
\label{sec:parametros}

Los parámetros de la EAM se fijaron en:
$\iota = 0$, $\kappa = 0$, $\sigma = 0.1$, y tolerancia de reconocimiento
$\xi = 0$ para el dominio textual y $\xi = 2$ para el directorio visual.
Estos valores siguen el modelo W-EAM y su extensión hetero-asociativa, descritos en el
marco teórico.
La propia línea W-EAM ha mostrado que estos hiperparámetros admiten optimización
automática multi-objetivo \cite{lopez2025hyperopt}; aquí se fijan por barrido directo,
cuya justificación experimental se presenta en la Sección~\ref{sec:param-iota-kappa}
(barrido $\iota \times \kappa$) y en la Sección~\ref{sec:hemisferio-visual}
(barrido de $\xi$ visual).
En síntesis: $\iota > 0$ elimina conocimiento escaso pero válido antes de eliminar
relaciones espurias, $\kappa$ no cambió ningún resultado en el rango probado, y
$\xi = 2$ en el directorio visual amplía la cobertura sin costo en especificidad.

% ---------------------------------------------------------------
\subsection{Auditoría de fidelidad}
\label{sec:auditoria}

Antes de producir los números citables se auditó que todas las decisiones del sistema
provengan de operaciones de las memorias y que ningún dato esté fabricado.
La auditoría encontró y corrigió un desvío real: el constructor de vectores de labels
admitía un \emph{fallback} sintético ($\pm 1$ por componente) para labels sin vector
fastText; un único label lo activaba y, además de esquivar el rechazo en consulta, sus
300 componentes de magnitud 1.0 inflaban el percentil 99 y \textbf{distorsionaban la
escala global de cuantización} de todos los labels ($S = 0.2276$ con el sintético
frente a $S = 0.1864$ real).
El fallback quedó deshabilitado (los labels sin vector se excluyen del llenado), las
memorias se rellenaron desde cero y toda la batería experimental se recorrió.
Los resultados de este reporte provienen exclusivamente de corridas posteriores a esa
auditoría; el impacto numérico de la contaminación resultó pequeño, pero los números
actuales son trazables a datos no fabricados.
La auditoría verificó además que el código de las memorias de Pineda \& Morales
permanece byte-idéntico al original y que la verdad de terreno del banco se usa
únicamente para calcular métricas, nunca dentro del circuito de decisión.

% ---------------------------------------------------------------
\subsection{Métricas de evaluación}
\label{sec:metricas}

Se emplearon las siguientes métricas:

\begin{itemize}
    \item \textbf{Precisión temprana} (\textit{early accuracy}): fracción de consultas para las que el agente ganador en la fase de activación temprana coincide con la clase real.
    \item \textbf{Precisión madura} (\textit{mature accuracy}): ídem tras la fase de recuperación madura, con lectura de directorio corregida por conteo (B1) o cruda.
    \item \textbf{Tasa de rechazo}: fracción de consultas para las que ningún agente supera el umbral de aceptación.
    \item \textbf{Fidelidad temprana-madura}: concordancia entre la clase ganadora en fase temprana y la clase ganadora en fase madura.
    \item \textbf{Entropía del directorio}: entropía de Shannon de la distribución de \textit{conteos} de los agentes especialistas, con máximo $\log_2(8) = 3$ bits.
    \item \textbf{Conteos por agente}: número de registros acumulados en el directorio transactivo por cada agente especialista.
    \item \textbf{Participación del ganador} (\textit{winner share}): fracción de victorias maduras capturadas por cada agente; el ideal balanceado es $1/8 = 12.5\%$.
    \item \textbf{Enrutamiento de test} (hemisferio visual): fracción de imágenes de prueba asignadas correctamente por el directorio visual, junto con la precisión sobre las imágenes aceptadas.
    \item \textbf{Evocación de dominio} (\textit{top-3 coseno}): fracción de imágenes para las que al menos una de las tres etiquetas más cercanas en el espacio fastText pertenece al vocabulario de la clase real (recuperación inversa imagen $\to$ palabras).
\end{itemize}

% ---------------------------------------------------------------
\subsection{Organización del experimento}
\label{sec:organizacion}

La evaluación se organiza en un estudio principal de calibración del directorio
(9 condiciones de lectura $\times$ $N \in \{50, 100, 200, 400\}$ consultas $\times$
5 semillas, sobre el banco de 411), cuatro secciones de caracterización bajo protocolo
único (protocolo completo, parámetros nativos $\iota \times \kappa$, formación del
directorio y capacidad de llenado, Secciones~\ref{sec:proto-completo}
a~\ref{sec:capacidad}), la evaluación del hemisferio visual
(Sección~\ref{sec:hemisferio-visual}) y la medición del costo computacional de la
reproducción completa (Sección~\ref{sec:costo}).
